\section{Discussion and Managerial Implications}
\label{sec:discussion}

The value of operational information depends on the managerial task it supports. The same operating state can inform which flight chains merit detailed review and whether a chain still offers a worthwhile local-recovery opportunity. These decisions compare different alternatives and operate at different decision boundaries. The value of an information representation should therefore be evaluated separately for screening and local recovery. This section discusses this distinction, decision-specific information value, and implications for airline operations control center (AOCC) decision support.

\subsection{Screening Priority versus Recoverable Opportunity}
\label{subsec:distinct_allocation}

Airline operations allocate scarce managerial resources in at least two ways: detailed-review attention across flight chains and recovery effort within an individual chain. Screening compares priorities across chains, whereas local recovery compares feasible targets within a chain. Baseline consequence severity is therefore distinct from recoverable value. High exposure may coexist with limited local improvement opportunities, while a chain with more moderate exposure may still admit a useful local adjustment. The two quantities address different managerial questions and should not be treated as interchangeable.

At the nominal review share $q_0=0.10$, delay-based screening replaces 37--55\% of the positions in the consequence-based reference shortlist. Adding passenger exposure reduces the screening loss, whereas adding downstream exposure produces results close to those of Delay alone. When the purpose of screening is to prioritize consequential disruptions, delay severity therefore does not fully substitute for consequence information.

Within the annual HJ reference cohorts, the preferred recovery target is positive under the Balanced objective for 82--92\% of cases. This result describes cases selected into the HJ reference cohorts; it does not establish the prevalence of positive recovery targets among all Stage-II-eligible chains, nor does it show how much recovery value a particular screening rule captures. Screening performance should be assessed against its intended purpose: consequence coverage for prioritizing consequential cases, and recovery-opportunity coverage when the goal is to identify cases with feasible improvement potential. The latter comparison cannot be estimated from the current held-out outputs.

\subsection{Decision-Specific Value of Operational Information}
\label{subsec:decision_specific_information_value}

More detailed information does not necessarily improve a decision. It has value for a particular task only when it changes the evaluation or ordering of the alternatives relevant to that task. The sufficiency of an information representation is therefore not an intrinsic property of the representation; it depends on the operating state, the objective, and the available alternatives.

In this study, HM and HJ produce nearly identical Stage-I shortlists: the reassigned share is 1.3--3.8\%, and $L^{\mathrm{att}}$ is 0.0006--0.0026. At Stage II, however, exact target agreement $A_0$ is 0.581--0.730, and $L^{\mathrm{rec}}$ is 0.0853--0.1247. Thus, close agreement in aggregate screening does not imply agreement on local-recovery targets.

The HP--HJ contrast is larger at both decision layers. It reassigns 84.2--88.7\% of shortlist positions, with $L^{\mathrm{att}}=0.4134$--$0.5291$; at Stage II, $A_0$ falls to 0.130--0.241 and $L^{\mathrm{rec}}$ ranges from 0.8963 to 0.9814. Within the tested objective and supported cohorts, compressing the state distribution to a point produces larger decision differences than replacing the joint distribution with weighted marginals. These are paired contrasts for the tested configurations, not a universal ranking of information representations. Stage-I and Stage-II losses also use distinct reference quantities and are not interpreted on a common scale.

The component audit offers a model-specific explanation for when joint dependence affects recovery evaluation. Under the weighted-Marginal coupling, Flight and Resource domain scores match their Joint counterparts to numerical precision for all audited candidates. Passenger domain scores differ for 39--66\% of candidates through threshold-event exposure. In the current consequence construction, joint dependence therefore affects evaluation primarily through nonseparable passenger-threshold consequences. This interpretation is specific to the modeled objective and does not establish a general mechanism for other consequence definitions or operating systems.

Objective weights further shape the decision consequences of information differences. Passenger emphasis yields the lowest HM--HJ target agreement in all six years, while Resource emphasis yields the lowest activation share and the largest zero-value share. Across the 24 year--objective settings, the median absolute target change is zero, even though the mean ranges from 1.8 to 4.7 minutes and the maximum from 30 to 45 minutes. Reporting only mean target change would therefore conceal the coexistence of many unchanged targets and a smaller number of larger changes. Objective weights, consequence structure, and information representation should be considered separately; a larger weight on one domain does not by itself imply greater sensitivity to joint dependence.

\subsection{Implications for AOCC Decision Support}
\label{subsec:managerial_implications}

At initial screening, delay signals can provide an interpretable basis for comparing many flights, but the current results do not support treating delay severity as a sufficient replacement for consequence priority. Passenger exposure improves the screening result in the tested comparisons, whereas downstream exposure does not show a comparable incremental gain. AOCCs may use consequence information to supplement delay-based ordering when the stated objective is to prioritize high-consequence cases, but these results do not validate a directly deployable screening policy.

Once a chain enters detailed diagnosis, the focus shifts from baseline exposure to remaining recovery opportunity, feasible local targets, and their modeled trade-offs. The results show that an information representation that preserves the screening order may still change the selected recovery target. Decision support for local recovery should therefore assess whether the representation distinguishes the feasible alternatives, rather than relying only on aggregate ranking agreement or descriptive completeness.

These findings support a principle of decision-matched information provision: screening emphasizes cross-flight consequence comparisons, while detailed diagnosis emphasizes remaining recovery opportunities and the evaluation of local targets. Cases near a screening cutoff or with closely evaluated recovery alternatives could be candidates for more detailed analysis. However, this study does not establish an optimal sequence for acquiring information or show that such sequencing reduces AOCC workload or improves executed recovery outcomes. Decision-matched information provision is therefore a candidate design principle for future decision-support systems, not a validated operating policy.

\subsection{Applicability and Future Research}
\label{subsec:boundary_conditions}

The evaluation uses publicly available U.S. airline operating records from 2017 to 2022 and compares decision relationships across years and PRE/TURN stages. These cross-year differences describe the observed operating environments and support populations; they do not identify the causal effect of any external shock. The experiments use prespecified canonical decision epochs and therefore do not describe how information and recovery opportunities co-evolve across successive rolling epochs within a flight chain.

Public operating records support reconstruction of flight execution, aircraft continuity, and several historical consequence-exposure measures. They do not fully capture carrier-internal information on crew legality, gate and ground-handling competition, maintenance restrictions, reserve aircraft, real-time passenger reaccommodation, or operational judgments. Such information may affect both consequence exposure and recovery feasibility. Access to carrier-internal operating and decision records would allow future research to assess the incremental decision value of these signals under more complete operating constraints.

The local-recovery decision is represented as an effective target constrained by modeled time headroom; the study does not translate that target into an executable resource plan. Incorporating aircraft swaps, crew adjustments, gate coordination, and shared ground resources would introduce more explicit cross-flight and cross-resource dependencies and could change the role of joint-state information in feasibility assessment. Future research could also evaluate recovery-opportunity coverage and the priority--recoverability relationship over the full Stage-II-eligible population; those quantities are not identified by the current held-out outputs.

Finally, the multidimensional consequence criterion uses transparent normalization and normative domain weights. Carrier-specific trade-offs among operational reliability, passenger service, and resource expenditure may depend on network structure, actual costs, service commitments, and operating policies. Linking the framework to carrier costs, managerial preferences, and executed recovery actions would help assess its applicability to specific organizations and test whether decision-matched information provision improves operational outcomes.
